% New standalone increment. Prior sources and GitHub are read-only.
\documentclass[11pt]{article}
\usepackage[T1]{fontenc}\usepackage[utf8]{inputenc}\usepackage{lmodern}
\usepackage[margin=.88in]{geometry}
\usepackage{amsmath,amssymb,amsthm,mathtools,microtype,array,longtable,booktabs,xurl}
\usepackage[hidelinks]{hyperref}
\setlength{\emergencystretch}{3em}\allowdisplaybreaks[2]
\numberwithin{equation}{section}
\newtheorem{theorem}{Theorem}[section]
\newtheorem{lemma}[theorem]{Lemma}
\newtheorem{proposition}[theorem]{Proposition}
\newtheorem{corollary}[theorem]{Corollary}
\theoremstyle{remark}\newtheorem{remark}[theorem]{Remark}
\newcommand{\E}{\mathbb E}\newcommand{\dd}{\,d}
\newcommand{\one}{\mathbf1}\newcommand{\eps}{\varepsilon}
\newcommand{\KK}{\mathcal K}\newcommand{\BB}{\mathcal B}
\newcommand{\CC}{\mathcal C}\newcommand{\Gp}{{G_+}}
\newcommand{\EE}{\mathfrak E}
\title{Probe-sector transverse energy and entry precision\\
in the initialized deep regime}
\author{Standalone analytical increment for independent review}
\date{Version 1 --- 2 November 2026}
\begin{document}\maketitle
\begin{abstract}
The lower-half input energy of every fixed cohort is bounded by
its initial energy plus \(Cq^2(1+\log(1/q))^2\) times its maximum
mass. The proof uses matrix values, all-label alignment and the
closed lower-side Gaussian gate; no receiver derivative or
ancestry sign is needed. A backward comparison from the exact
radial equation pays all pre-entry cohort masses using the
already proved passage budgets. Consequently every outside
probe moment left open by CWI is \(o(q^2)\) throughout the late
weak passage. The full initialized error and source rates match
the selected leading orbit at every fixed sufficiently small
logistic level, uniformly on compact sets of such levels.
A separate conditional estimate identifies the bounds needed
on a later fixed window; no continuation beyond \(t_b\) is proved.
\end{abstract}

\section{Conventions, interval and hypothesis inventory}
\label{an07psev1:inventory}
Keep the original balanced, untied flow with
\[
 P_{\rm mix}=\tfrac12(P_1+P_2),\quad
 P_1=N(e_1,q^2I_2),\quad P_2=N(\rho e_2,q^2I_2),\quad
 \rho\in[13/20,7/10],\quad S_0=q^{10}.
 \]
Time is normalized: \(t=\tau=\mu_1^2t_{\rm phys}\).
Write \(\lambda=\rho^2,\ Q=\lambda+q^2,\ \ell=\log(1/q)\),
and \(L=1+\ell\). All estimates are for sufficiently small
\(q>0\), uniformly on the ratio box after the unchanged
WPT/RB and WCS constant choices.
Let \(U=\sqrt m\,a_\theta,\ W=\sqrt m\,a_\psi\), with
\(a_\theta=(\cos\theta,\sin\theta)\) and
\(\dd\lambda_0=\dd b/(2\pi)\). The fixed partition is
\(\KK,\Gp,\Sigma,\BB\), with selected clock \(S=\mathcal C_q\).
Fixed cohorts may depend on \(q\); they do not depend on time.
The strong complement is \(\CC\subset\BB\).

The proved times are
\[
 t_m=\left(\frac4{1-\lambda}+\frac12\right)\ell,\quad
 t_0=10\ell-\log A_q+o(1),\quad
 t_b<t_0+20\ell .
 \]
Here \(0<c_A\le A_q\le C_A\) is MOMC's entry constant.
Write \(u=19/2-4/(1-\lambda)\ge169/102\) for the early
mass exponent.
In particular \(t_b\le31L\) for small \(q\).
Let \(t_s=\inf\{t\ge t_0:H(t)=q^{1/8}\}\), as in WCS.
It lies before \(t_b\), since the proved clock has
\(H_0\asymp A_q^{-\lambda}q^{10(1-\lambda)}\) and
\(H(t_b)\asymp c_b>0\).
Fix a probe range \(0\le\zeta\le R_\zeta\), with \(R_\zeta\)
independent of \(q\), and require \(qR_\zeta\le\pi/2\).
This explicit sign restriction includes WIT's positive offsets.

For a fixed cohort \(T\) define
\[
 \mu_T=\int_Tm\dd\lambda_0,\qquad
 \EE_T=\int_T(U_2)_-^2\dd\lambda_0,\qquad
 \mathcal Q_T(\zeta)=
 \int_Tm[\sin(q\zeta-\theta)]_+^2\dd\lambda_0 .
 \]
The last quantity is exactly CWI's probe moment, not a mass
majorant. The symbol \(e_f=X-f(X)\) denotes the residual.
Matrix conventions are
\begin{equation}
 R^{(j)}_\theta=\tfrac12\E_{P_j}
       [e_f(X)X^T\one_{\{U\cdot X>0\}}],\quad
 R=R^{(1)}+R^{(2)},\quad U'=R^TW,\quad W'=RU .
 \label{an07psev1:matrixconvention}
\end{equation}
The mixture half is included in each \(R^{(j)}\).

\small
\begin{longtable}{>{\raggedright\arraybackslash}p{.22\textwidth}
 >{\raggedright\arraybackslash}p{.39\textwidth}
 >{\raggedright\arraybackslash}p{.31\textwidth}}
\toprule Source & Hypotheses needed here & Interval and supply\\\midrule\endhead
P \texttt{pa:exactflow}, loss and balance
& Initialized positive balanced radii; Gaussian model and normalized
 time; regular characteristic flow and finite moments.
& All existing times through \(t_b\). Exact identities;
 loss monotonicity is displayed below. No angular sign required.\\
P \texttt{pa:earlybounds}
& Aligned initialization, \(S_0=q^{10}\); for its angle bound,
 \(\overline S-S_0<\pi\).
& \([0,t_m]\): \(\overline S(t_m)=O(q^u)\),
 \(u\ge169/102>1\), verifies the range condition.\\
MOMC \texttt{main, massclaim, chartclaim, outsideclaim}
& Original initialized deep regime, small \(q\); fixed right-half
 ancestry \(\mathcal R\), positive radial weights, first right-half
 mass \(1/2\).
& Unconditional S1 result on \([t_m,t_0]\):
 \(0<M_{\mathcal R}\le1/2\), chart \(79/51\),
 \(\mu_{\mathcal R^c}\le Cq^2\ell^2M_{\mathcal R}\).
 Its integrated mass consequence is proved below.\\
Q3SB \texttt{pre/crosspre}
& Same initialized model and the preceding early/MOMC results;
 arbitrary receiver, no pre-entry response sign.
& \([0,t_0]\): \(\|e_1\|_{P_2}+\|e_2\|_{P_1}\le C_{\rm pre}q\).
 No strip-entry assumption.\\
Q2EC \texttt{fullgates}, RB \texttt{pretrace}
& Strong donors in their compact chart, \(M_{\mathcal R}\le1/2\),
 left mass charged absolutely; arbitrary receiver gates.
& \([t_m,t_0]\). The proof-level intermediate trace estimate
 is restated and justified below, then integrated at every time.\\
RB \texttt{completed, conclusions, massconclusions, clock}
& Initialized deep regime; fixed WPT/RB order and sufficiently
 small \(q\). Its lower response barrier is proved, not retained.
& \([t_0,t_b]\): \(0<c\le c_b\), duration, residual, mass,
 core chart and clock conclusions. No new stopped guard is imposed.\\
RB \texttt{postalignment}
& Completed passage supplies its cross-residual and signed-gate
 budgets; continuous lifts and initial \(O(q)\) alignment.
& \([t_0,t_b]\): \(|\theta-\psi|\le C_{\rm al}qL\), all labels.
 No outside sign assumption.\\
S2CP energy/mass, PISB/LSB/GESG budgets
& Their stopped hypotheses, including \(0\le c\le c_b\), entry,
 geometry, field and strip budgets, were discharged by WPT/RB.
& Used only as RB's proved passage conclusions. Their extension
 to \([0,t_0]\) follows from the new radial comparison, not
 from applying a stopped theorem before its start.\\
CWI \texttt{split, equivalence}
& Balanced exact output; fixed \(T=\Sigma\cup\BB\);
 inactive \(\Gp\); \(|f_\KK|=O(q)\); alignment \(\delta_q\);
 \(\mu_TL^2\to0\) for the equivalence.
& Originally at \(t_n\). These hypotheses are verified directly
 on \([t_s,t_b]\) below, so its algebra applies there.\\
CWI \texttt{ratecharge}
& Sourcewise bounded diagonal and \(O(q)\) cross entries;
 alignment \(CqL\); fixed cohort, bounded offset, selector/a.e.
 convention at the probe gate.
& Passage matrix bounds supply all hypotheses on \([t_s,t_b]\).
 Only values, not receiver derivatives, are used.\\
CWI \texttt{corestatic, corerates}
& Reached core chart; physical core mass and first moments close
 to the selected state; same \(\rho\), offset and phase.
 Rates also require a fixed positive core probe gap, sourcewise
 CFA reduction and its pointwise late remainder, SRR moments.
& At \(t_n\) in CWI. After the uniform-level lemma below,
 all hypotheses hold at \(t_{n'}\), \(n'\) in a fixed compact
 subset of \((0,n]\). No later-window comparison is inferred.\\
WCS/SRR/MP/WBR and TAILREP \texttt{selection}
& Exact auxiliary logistic, \(0<N_0<n'<1\); maintained original
 return and local strong box; positive measures and core moment;
 finite \(L^2(\eta_s)\), fixed incoming law; both error inequalities,
 nonnegative \(\ell_B\), \(\int\ell_B\le1\), common pointwise
 \(d_0\); small gain and initial strict guard.
& \([t_s,t_{n'}]\subset[t_s,t_n]\subset[t_s,t_e]\).
 WCS discharges the field/return inputs. The reduced endpoint
 small-gain and uniformity checks are given in Section 5.\\
\bottomrule
\end{longtable}\normalsize

\section{Matrix values, alignment and pre-entry cohort masses}
\label{an07psev1:inputs}
\subsection{Loss and cross residuals on their correct intervals}
For the normalized training loss
\(\mathcal L=\tfrac12\E_{\rm mix}|X-f(X)|^2\), the exact
gradient-flow identity is
\[
 \mathcal L'=-\int(|U'|^2+|W'|^2)\dd\lambda_0\le0.
 \]
At aligned isotropic initialization \(f_0(X)=q^{10}X/4\).
For \(q\le1/2,\ \rho\le1\), this gives \(\mathcal L(0)\le1\);
therefore \(\|e_f\|_{L^2(P_j)}\le2\) for \(j=1,2\).
These bounds hold at every time in the proved trajectory.

Set \(C_\times=\max(C_{\rm pre},M_{\rm res})\).
Q3SB \texttt{crosspre} and RB \texttt{conclusions} give
\begin{equation}
 \|e_2\|_{P_1}+\|e_1\|_{P_2}\le C_\times q
                 \qquad(0\le t\le t_b).
 \label{an07psev1:cross}
\end{equation}
The pre-entry source's two regimes are essential:
on \([0,t_m]\), \(S\le q^{10}e^{(1+2q^2)t}=O(q^u)=o(q)\)
and \(|f(X)|\le S|X|\); on \([t_m,t_0]\),
\[
 \|(f_{\mathcal R})_1\|_{P_2}
 +\|(f_{\mathcal R})_2\|_{P_1}\le CqM_{\mathcal R},
 \quad
 \|(f_{\mathcal R^c})_i\|_{P_j}\le C\mu_{\mathcal R^c}
 \le Cq^2\ell^2M_{\mathcal R}.
 \]
Thus both cross target coordinates, whose norms equal \(q\),
have \(O(q)\) residuals. This is the content of Q3SB's proved
lemma, not an application of LCRC before \(t_0\).
In particular \(\|f_2\|_{P_1}\le(1+C_\times)q\);
the \(P_1\) second-output term cannot have an order-one price.

\begin{lemma}[Uniform matrix values, with the lower gate retained]
\label{an07psev1:matrix}
For all receivers and \(0\le t\le t_b\),
\begin{equation}
 \begin{gathered}
 |R_{12}|,\,|R_{21}|\le C_{12}q,\qquad
 |R_{12}^{(j)}|,\,|R_{21}^{(j)}|\le C_{12}q\quad(j=1,2),\\
 |R_{22}^{(1)}|\le C_1q^2,\qquad
 C_{12}=1+C_\times,\quad C_1=C_\times/2.
 \end{gathered}
 \label{an07psev1:matrixbounds}
\end{equation}
All entries of each source matrix are bounded by a model
constant. If \(U_2<0\), put \(s=-\sin\theta>0\). Then
\begin{equation}
 \mathbb P_{P_2}(U\cdot X>0)=\Phi(-\rho s/q),\qquad
 |R_{22}^{(2)}|\le C_G e^{-k_Gs^2/q^2},
 \quad C_G=2,\quad k_G=\frac{169}{3200}.
 \label{an07psev1:gate}
\end{equation}
\end{lemma}
\begin{proof}
For \(R_{12}^{(1)}\), use \(\|e_1\|_{P_1}\le2\) and
\(\|X_2\|_{P_1}=q\). For \(R_{12}^{(2)}\), use the cross
residual and \(\|X_2\|_{P_2}=\sqrt Q\le1\).
For \(R_{21}\) interchange the clusters, using
\(\|X_1\|_{P_1}\le2\). This proves the first two bounds
with the stated slack. Cauchy--Schwarz gives
\(|R_{22}^{(1)}|\le\frac12\|e_2\|_{P_1}q\).
The loss bounds give the remaining bounded entries.

The exact receiver-gate probability follows from
\(a_\theta\cdot X\sim N(\rho\sin\theta,q^2)\) under \(P_2\).
For \(A=\{U\cdot X>0\}\), Hölder and Cauchy--Schwarz give
\[
 |R_{22}^{(2)}|
 \le\tfrac12\|e_2\|_{P_2}
          (\E_{P_2}X_2^4)^{1/4}\mathbb P(A)^{1/4}.
 \]
Here \(\E X_2^4=\rho^4+6\rho^2q^2+3q^4\le43/16\)
for \(\rho\le1,\ q\le1/2\), so its fourth root is less than
two. The Gaussian Chernoff bound
\(\Phi(-v)\le e^{-v^2/2}\) for \(v\ge0\) proves
\eqref{an07psev1:gate}, including the mixture half.
No independence of \(e_2\) and the gate, or integration by
parts across that gate, was asserted.
\end{proof}

\subsection{Alignment before entry, including intermediate times}
\begin{lemma}[All-label alignment over all three regimes]
\label{an07psev1:alignment}
There is a fixed \(C_{\rm a}\) such that
\[
 |\theta-\psi|\le
 \begin{cases}
 Cq^u,&0\le t\le t_m,\\
 Cq,&t_m\le t\le t_0,\\
 C_{\rm a}qL,&t_0\le t\le t_b .
 \end{cases}
 \qquad
 |W_2-U_2|\le C_{\rm a}qL\sqrt m .
 \label{an07psev1:alignmentbounds}
 \]
All initial labels and closed sector endpoints are included.
\end{lemma}
\begin{proof}
P \texttt{earlybounds} applies through \(t_m\), since
\(\overline S(t_m)=O(q^u)<\pi\), proving the first line.
For the intermediate line we restate the calculation in
RB \texttt{entryalignment}; an endpoint estimate alone would
not suffice.
Under MOMC's theorem, let
\(T_{\mathcal R}=\int_{\mathcal R}m a_\psi a_\theta^T\dd\lambda_0\).
Q2EC \texttt{fullgates} gives, for every receiver,
\[
 R=\tfrac12[(I-T_{\mathcal R})A_1+A_2-B_2]+\mathcal R,\quad
 \|\mathcal R\|\le C\mu_{\mathcal R^c}
                         +CM_{\mathcal R}e^{-c/q^2}.
 \]
The target \(A_j\) are positive semidefinite and symmetric.
The core diagonal coefficients are \(M_{\mathcal R}+O(q^2M_{\mathcal R})\)
and \(O(q^2M_{\mathcal R})\); its cross coefficients are
\(O(qM_{\mathcal R})\). The corresponding \(B_2\) diagonal
entries are \(O(q^2M_{\mathcal R})\) and its cross entries
are \(O(qM_{\mathcal R})\).
Since \(M_{\mathcal R}\le1/2\), the diagonal coefficients
of \(I-T_{\mathcal R}\) are positive for small \(q\).
Also \(|(A_1)_{12}|\le Cq\). Thus
\begin{align*}
 (\operatorname{tr}R)_-
 &\le C(q^2M_{\mathcal R}+\mu_{\mathcal R^c}
                             +M_{\mathcal R}e^{-c/q^2}),\\
 |R_{12}-R_{21}|
 &\le C(qM_{\mathcal R}+\mu_{\mathcal R^c}
                             +M_{\mathcal R}e^{-c/q^2}).
\end{align*}
These are exactly RB's intermediate trace bounds, now with
their needed donor and receiver hypotheses explicit.

For clarity, MOMC's mass law already yields
\[
 \int_{t_m}^{t_0}M_{\mathcal R}
 =\log\frac{1-M_0}{1-1/2}
   -\int_{t_m}^{t_0}
       \frac{M_{\mathcal R}\delta_M}{1-M_{\mathcal R}}
 \le\log2+\int|\delta_M|=\log2+o(1).
 \]
Here \(M_{\mathcal R}>0,\ M_{\mathcal R}\le1/2\) and
\(\int|\delta_M|=O(q^2\ell^2)\), all from its theorem.
The negative trace integral is therefore \(O(q^2\ell^2)\)
and the antisymmetric integral is \(O(q+q^2\ell^2)=O(q)\).
The exact difference equation, for \(\Delta=\theta-\psi\), is
\[
 \Delta'=-\operatorname{tr}R\sin\Delta
                         +(R_{12}-R_{21})\cos\Delta.
 \]
While \(|\Delta|<1\), its upper Dini derivative is at most
\((\operatorname{tr}R)_-|\Delta|+|R_{12}-R_{21}|\).
Gronwall from \(t_m\) gives \(Cq\) uniformly at every
intermediate time, strictly closing that temporary guard.
The post-entry line is RB \texttt{postalignment} on its
completed interval. Finally sine is 1-Lipschitz and balance
gives the stated coordinate difference.
\end{proof}

\begin{lemma}[Pre-entry mass comparison for every fixed cohort]
\label{an07psev1:backmass}
For sufficiently small \(q\), every label and every \(0\le t\le t_0\)
satisfies
\[
 m(t,b)\le e^5m(t_0,b),\qquad
 \sup_{0\le t\le t_0}\mu_T(t)\le e^5\mu_T(t_0)
 \quad\hbox{for every fixed }T .
 \label{an07psev1:backmassbound}
 \]
No backward clock premise or future probe bound is used.
\end{lemma}
\begin{proof}
P's exact decomposition \(R=A_\theta-C_f\) has
\(A_\theta\succeq0,\ \|A_\theta\|\le a_*=\frac12+q^2\)
and \(\|C_f\|\le a_*S(t)\), \(S=\int m\).
The last inequality follows from the feature representation
and Gaussian second-moment Cauchy--Schwarz, for arbitrary
receiver gates. The symmetric quadratic identity gives
\[
 (\log m)'=2a_\psi^TRa_\theta
          \ge-a_*|\theta-\psi|^2-2a_*S.
 \]
The early envelope has integral \(O(q^u)\) through \(t_m\).
The preceding mass calculation and MOMC's outside bound give
\(\int_{t_m}^{t_0}S\le\log2+o(1)\), hence
\(\int_0^{t_0}S\le2\) for small \(q\).
Lemma \ref{an07psev1:alignment} gives
\(\int_0^{t_0}|\theta-\psi|^2\le Cq^2L\le1\), uniformly
in the label. Since \(a_*\le1\), integration from \(t\) to
\(t_0\) gives \(\log(m(t_0)/m(t))\ge-5\).
Integrate the labelwise bound over \(T\).
\end{proof}

\section{Static probe inequality and the energy theorem}
\label{an07psev1:energysection}
\begin{lemma}[Static inequality with the exact sign range]
\label{an07psev1:static}
For every balanced state and \(0\le\zeta\le R_\zeta\),
\(qR_\zeta\le\pi/2\),
\begin{align}
 (U\cdot\xi_q)_+
 &\le |U_1|\sin(q\zeta)+(U_2)_-\cos(q\zeta),\notag\\
 \mathcal Q_T(\zeta)&\le2q^2\zeta^2\mu_T+2\EE_T .
 \label{an07psev1:staticbound}
\end{align}
\end{lemma}
\begin{proof}
Both trigonometric coefficients on the right are nonnegative.
Use \(U\cdot\xi_q=U_1\sin(q\zeta)-U_2\cos(q\zeta)\),
then square, apply \((a+b)^2\le2a^2+2b^2\),
\(\sin(q\zeta)\le q\zeta\), and \(U_1^2\le m\).
For \(\theta=0\) the first inequality is equality;
for \(\theta=\pi\) its left side is zero; for
\(\theta=3\pi/2\) it is again equality, with value
\(\sqrt m\cos(q\zeta)\). At \(\zeta=0\), the first two
angles are gate endpoints. The classification uses current
angles, not initial ancestry. No version with unrestricted
signed \(\zeta\) or negative cosine is used.
\end{proof}

\begin{lemma}[Exact negative-input energy identity]
\label{an07psev1:identity}
Let \(h=(U_2)_-\). Along each characteristic,
\begin{equation}
 (h^2)'=-2h(W_1R_{12}+W_2R_{22}),\qquad
 R_{22}=R_{22}^{(1)}+R_{22}^{(2)} .
 \label{an07psev1:energyidentity}
\end{equation}
It holds almost everywhere, with both sides zero from the
chain rule at \(U_2=0\). The integrated identity holds for
every fixed cohort. There is no moving-boundary flux.
\end{lemma}
\begin{proof}
The function \(v\mapsto(v_-)^2\) is continuously differentiable
with derivative \(-2v_-\). Apply the chain rule to
\(U_2'=W_1R_{12}+W_2R_{22}\).
At each fixed \(q\), bounded residual matrices give an
exponential finite-time radial bound, justifying
integration over the finite label law. No differentiation
of the indicator of a time-varying cohort is involved.
\end{proof}

\begin{theorem}[Lower-half energy, with \(j=2\)]
\label{an07psev1:energy}
For every fixed cohort \(T\) and every \(0\le t\le t_b\),
\begin{equation}
 \EE_T(t)\le e^{31}\EE_T(0)
       +31e^{31}C_Eq^2L^2
                       \sup_{0\le v\le t}\mu_T(v),
 \qquad \EE_T(0)\le q^{10},
 \label{an07psev1:energybound}
\end{equation}
where, with the constants of Lemmas \ref{an07psev1:matrix}
and \ref{an07psev1:alignment}, one may take
\begin{equation}
 C_E=C_{12}^2+2C_1+
              \frac{2C_G}{e k_G}
              +\frac{2C_G C_{\rm a}}{\sqrt{2e k_G}}.
 \label{an07psev1:CE}
\end{equation}
All constants are fixed independently of \(q\) and the cohort.
\end{theorem}
\begin{proof}
On \(U_2<0\) put \(s=h/\sqrt m=-\sin\theta\).
Balance gives \(|W_1|\le\sqrt m\) and
\(|W_2|\le h+\delta\sqrt m\), where
\(\delta=|\theta-\psi|\le C_{\rm a}qL\).
The cross term in \eqref{an07psev1:energyidentity} obeys
\[
 2h|W_1R_{12}|
 \le L^{-1}h^2+C_{12}^2q^2Lm .
 \]
The \(P_1\) diagonal term is at most \(2C_1q^2m\), using
\(h,|W_2|\le\sqrt m\), without an output sign assumption.
The \(P_2\) term is at most
\[
 2C_Gm(s^2+\delta s)e^{-k_Gs^2/q^2}
 \le2C_Gm\left(\frac{q^2}{e k_G}
                 +\frac{\delta q}{\sqrt{2e k_G}}\right).
 \]
The last inequality follows by maximizing
\(v^2e^{-k_Gv^2}\) and \(ve^{-k_Gv^2}\) for \(v\ge0\);
their maxima are \(1/(e k_G)\) and \(1/\sqrt{2e k_G}\).
This retains the Gaussian factor rather than treating
the \(h^2R_{22}^{(2)}\) term as an order-one growth rate.
Thus, including \(U_2\ge0\) by the zero derivative there,
\begin{equation}
 (\EE_T)'\le L^{-1}\EE_T+C_Eq^2L\mu_T
                       \quad\hbox{a.e. on }[0,t_b].
 \label{an07psev1:energydiff}
\end{equation}
Integrate and use \(t_b\le31L\).
Initially \(m=q^{10}\) and \(\lambda_0\) is a probability,
so \(\EE_T(0)\le q^{10}\). In fact \(\EE_\Sigma(0)=0\),
because its initial input angles lie in Q2.
\end{proof}

\paragraph{Explicit layer check.}
The suggested split is valid with
\(K_{\rm lay}^2=4/k_G=12800/169\).
On \(s\le K_{\rm lay}q\sqrt\ell\) one has
\(h^2\le K_{\rm lay}^2q^2\ell m\); on its other side
\(|R_{22}^{(2)}|\le2q^4\).
The separating level is below one for small \(q\).
Both sides, and \(s=0\), are covered by the proof's stronger
Gaussian maxima. This is why no third logarithmic power
is needed. Output misalignment was charged even when
\(W_2\) has the opposite sign from \(U_2\).

\section{All cohort masses and the late-passage probe budget}
\label{an07psev1:budgets}
Let \(\gamma_*:=289169/924000>3/10\) and
\[
 \Gamma_\nu=1-5\lambda+\frac{2\lambda}{1-\lambda}
                   -\nu[10(1-\lambda)-2]\ge\gamma_*,
 \qquad 0<\nu\le1/100,\quad q^2/\lambda\le\nu .
 \]
These are the unchanged S2CP choices. Its energy conclusion,
inherited by RB, and \(J_\Sigma\le J_G\le q^2K_Y^2\), give
\[
 \sup_{[t_0,t_b]}\mu_\Sigma
 \le C_\Sigma(q^5+q^{\Gamma_\nu}E_g)+q^2K_Y^2 ,
 \]
enlarging only a fixed coefficient; \(H\le E_g\) for small
\(q\) follows from RB with \(E_g=2c_b\).
Each mass estimate below is a proved RB passage conclusion.
Lemma \ref{an07psev1:backmass} multiplies its right side
by at most \(e^5\) to cover the entire interval \([0,t_b]\).
No source is applied before its proved time range.
\small
\begin{longtable}{>{\raggedright\arraybackslash}p{.24\textwidth}
 >{\raggedright\arraybackslash}p{.65\textwidth}}
\toprule Fixed cohort & Bound for \(\mu_T^*=\sup_{[0,t_b]}\mu_T\)\\\midrule\endhead
\(\Sigma\) & \(C[q^5+q^{\Gamma_\nu}E_g+q^2K_Y^2]\);
 S2CP/RB and the backward comparison.\\
Strong complement \(\CC\) & \(Cq^{1+\Delta_\nu}\),
 \(\Delta_\nu\ge17141311/140024000>3/25\);
 S2CP/RB and the backward comparison.\\
\(\pi\) strip & \(C[q^{1211/800}(1+\sqrt{\log\ell})
                          +q^{11}\sqrt\ell]\);
 PISB as inherited by RB.\\
\(J_3\) inner zone & \(Cq^{4769/4000}(1+\sqrt{\log\ell})\);
 LSB as inherited by RB.\\
Q3 interior & \(Cq^{10}\); GESG as inherited by RB.\\
Frozen zones & \(Cq^{11}\sqrt\ell\); PISB as inherited by RB.
 They are included in whole-strip totals and not counted twice.\\
Entire \(\BB\), or any remaining subcohort & \(e^5q^{11/10}\)
 for small \(q\), by RB's aggregate bound and the backward
 comparison. The fixed ancestry decomposition is exhaustive
 up to the already assigned zero-measure endpoints.\\\bottomrule
\end{longtable}\normalsize
Every listed mass times \(L^2\) tends to zero. For
\(T=\Sigma\cup\BB\) in particular,
\begin{equation}
 \mu_T^*\le C(q^5+q^{\Gamma_\nu}+q^2+q^{11/10})
                 \le Cq^{\gamma_*}.
 \label{an07psev1:totalmass}
\end{equation}
The last comparison uses \(0<q<1\) and
\(\gamma_*<11/10<2<5,\ \Gamma_\nu\ge\gamma_*\);
no smaller error power has silently absorbed a larger one.
Combining \eqref{an07psev1:staticbound} and
\eqref{an07psev1:energybound} proves the stronger all-time bound
\begin{equation}
 \sup_{\substack{0\le t\le t_b\\0\le\zeta\le R_\zeta}}
 \mathcal Q_{\Sigma\cup\BB}(t,\zeta)
 \le C_{R_\zeta}q^2\epsilon_{\rm bg}(q)=o(q^2),
 \quad
 \epsilon_{\rm bg}(q):=q^8+L^2q^{\gamma_*}.
 \label{an07psev1:Qpaid}
\end{equation}
This pays current-angle lower-side slices and all transit
labels, without assigning them a new sign or chart.

\begin{lemma}[Primary inactivity over the whole late interval]
\label{an07psev1:inactive}
All labels in \(\Gp\) are probe-inactive for
\(t_s\le t\le t_b\), uniformly on any fixed bounded offset range.
\end{lemma}
\begin{proof}
RB's completed RCET transport has \(U_2>0\) and
\[
 q|U_1|/U_2\le C(C_0e^{-b_*(t-t_0)/2}+q^2),\qquad
 b_*=Q(1-c_b)/2 .
 \]
The exact subinterval clock, with \(I'\le Q\), gives
\[
 t_s-t_0\ge Q^{-1}[(k-1/8)\ell+O(1)],
 \quad k=10(1-\lambda).
 \]
Since \(c_b\le1/4\) and \(k\ge51/10\),
\[
 e^{-b_*(t_s-t_0)}
 \le Cq^{(1-c_b)(k-1/8)/2}\le Cq^{597/320}.
 \]
This repeats SRR's late-time algebra on its now stated
scope, using only RB/RCET conclusions valid to \(t_b\).
For \(t\ge t_s\), \(q|U_1|/U_2=o(1)\); hence
\[
 U\cdot\xi_q/U_2
 \le R_\zeta q|U_1|/U_2-\cos(q\zeta)<0 .
 \]
The positive coordinate and regularized ratio include
the closed boundary label \(b=\pi\).
\end{proof}

\begin{theorem}[Full background precision within the weak passage]
\label{an07psev1:lateprecision}
Uniformly for \(t\in[t_s,t_b]\), \(0\le\zeta\le R_\zeta\),
\begin{align}
 |E_{\xi_q}(t)-E_{\xi_q}^{\KK}(t)|
 &\le C_{R_\zeta}q^2\epsilon_{\rm bg}(q),\notag\\
 |\xi_q\cdot\dot f_{\Sigma\cup\BB,p}|
 &\le C_{R_\zeta}q^2\epsilon_{\rm bg}(q),\notag\\
 |D_p^\tau-D_p^{\KK,\rm self}|
 &\le C_{R_\zeta}q^2\epsilon_{\rm bg}(q),\qquad p=1,2 .
 \label{an07psev1:latebounds}
\end{align}
Here \(E^\KK=\tfrac12|\xi_q-f_\KK|^2\) and
\(D_p^{\KK,\rm self}=-(\xi_q-f_\KK)\cdot\dot f_{\KK,p}\);
the core velocities are driven by the full original field.
Use the exact selector source rates or a.e. derivatives.
No uniform pointwise claim at gate-hit atoms is made.
\end{theorem}
\begin{proof}
The preceding lemma makes \(f=f_\KK+f_T\), \(T=\Sigma\cup\BB\),
on this interval. RB supplies the core chart and bounded
core mass, so \(|f_\KK|\le Cq\).
CWI's exact split, with all its hypotheses now verified, gives
\[
 E-E^\KK=-\mathcal Q_T+\mathcal A_T
                      +f_\KK\cdot f_T+\tfrac12|f_T|^2,\quad
 |\mathcal A_T|\le CqL\sqrt{\mu_T\mathcal Q_T},\quad
 |f_T|\le\sqrt{\mu_T\mathcal Q_T}.
 \]
Insert \eqref{an07psev1:Qpaid} and
\(\mu_T\le Cq^{\gamma_*}\). In particular
\[
 L\sqrt{\mu_T\epsilon_{\rm bg}}
 \le C(q^4L\sqrt{\mu_T}+L^2\mu_T)
 \le C(q^8+L^2q^{\gamma_*}).
 \]
The last step is Young's inequality. It pays every cross
and misalignment term at the stated precision.

All sourcewise matrix hypotheses of CWI \texttt{ratecharge}
hold by Lemma \ref{an07psev1:matrix}. Its conclusion is
\[
 |\xi_q\cdot\dot f_{T,p}|
 \le C(\mathcal Q_T+q^2L^2\mu_T),\qquad
 |\dot f_{T,p}|
 \le C(\sqrt{\mu_T\mathcal Q_T}+qL\mu_T).
 \]
It applies to arbitrary receivers and so needs no outside
chart. At core receivers, bounded diagonal and \(O(q)\)
cross entries give \(|\theta'_p|+|\psi'_p|=O(q)\) and
\(|(\log m)'_p|=O(1)\). The exact output derivative therefore
has \(|\dot f_{\KK,p}|\le Cq\).
Finally expand exactly
\[
 D_p^\tau-D_p^{\KK,\rm self}
 =f_T\cdot\dot f_{\KK,p}
  -(\xi_q-f_\KK)\cdot\dot f_{T,p}
  +f_T\cdot\dot f_{T,p}.
 \]
The displayed bounds and the same Young estimate pay all
three terms. The invisible primary labels remain in the
training residual, even though their direct probe rates are zero.
\end{proof}
Thus on the late passage the full error has the instantaneous
core-chart expansion
\[
 q^{-2}(E_{\xi_q}-1/2)=\mathcal E_{\KK(t),\zeta}
                       +O(q^2+\epsilon_{\rm bg}).
 \]
This is not selection of that core throughout \([t_n,t_b]\).
Selected-orbit comparisons are asserted only at the levels below.

\section{Uniform entry comparison at fixed smaller levels}
\label{an07psev1:entrysection}
Retain WCS's same auxiliary logistic \(N(t_s)=N_0=q^{1/8}\),
and its fixed upper matching level
\[
 \begin{gathered}
 n=\frac12\min\left\{h,\frac14,
       \frac{\lambda_{\min}g_{\min}}
            {64C_sB_{\max}C_w}\right\},\\
 \lambda_{\min}=\frac{169}{400},\
 g_{\min}=\frac{1183}{41600},\
 C_s=980,\ C_w=20,\ 
 B_{\max}=e\,2^{1/(2\lambda_{\min})}.
 \end{gathered}
 \]
Fix a compact level set \(\mathcal N\subset(0,n]\), and
write \(n_-=\inf\mathcal N>0\).
For \(n'\in\mathcal N\), let \(t_{n'}\) be the hit of this
same logistic mass. For small \(q\), \(N_0<n_-\) and
\([t_s,t_{n'}]\subset[t_s,t_n]\subset[t_s,t_e]\).

\begin{lemma}[WCS at every such level, with uniform constants]
\label{an07psev1:levels}
Uniformly in \(n'\in\mathcal N\) and \(\rho\) in the ratio box,
\[
 E_s(t_{n'})+W_2(t_{n'})=O_{\mathcal N}(q^{1/20}).
 \label{an07psev1:levelselection}
 \]
WCS's other entry contracts and pointwise late bounds used
by CWI hold at these times with constants uniform on
\(\mathcal N\).
\end{lemma}
\begin{proof}
Use the already proved WCS strong and weak error inequalities
on the smaller interval, with the unchanged positive logistic
comparator, actual returned weak inputs and positive donor
measure, local strong box, incoming-law \(L^2\) norm,
\(\ell_B\ge0,\ \int\ell_B\le1\), and common pointwise
\(d_0=Cq^{1/20}\). WCS's initial estimates are
\(E_0=O(q^{1/20})\) and \(A_0=O(q^{1/16})\).
These are conclusions under the original-flow hypotheses,
not conditions added to this lemma.

For TAILREP's endpoint constants with \(n'\) in place of \(n\),
\[
 B_*\le B_{\max},\quad a_0=\lambda(1-n')/2,\quad
 m_0=\min(g/2,a_0/2),\quad
 a_0-m_0\ge\lambda_{\min}/8,\quad g-m_0\ge g_{\min}/2 .
 \]
Consequently
\[
 \Theta(n')\le
 \frac{16C_sB_{\max}C_wn'}{\lambda_{\min}g_{\min}}
 \le\frac18<\frac14 .
 \]
\(K_0=2[1+C_s(1+B_*)/(g-m_0)]\) is uniformly bounded,
and its initial strict guard
\(K_0(E_0+d_0+A_0\sqrt{n'})\le\delta/2\) holds uniformly
for small \(q\). The return, finite moments, regularity and
first-exit hypotheses of TAILREP are exactly the ones already
verified in WCS; restriction of the time interval preserves them.
Thus its two conclusions apply verbatim with endpoint \(n'\).
In particular
\[
 W_2(t_{n'})\le B_*\left[
 \frac{A_0}{\sqrt{n'}}+
 C_w(K_0+1)\left\{
 \frac{E_0e^{-m_0(t_{n'}-t_s)}}{a_0-m_0}
 +\frac{d_0}{a_0}+\frac{A_0\sqrt{n'}}{2a_0}\right\}\right].
 \]
All denominators are uniformly positive, including
\(\sqrt{n'}\ge\sqrt{n_-}\). Since \(q^{1/16}\le q^{1/20}\),
this and the strong-error conclusion prove the stated rate.
Also
\[
 t_{n'}-t_s=\lambda^{-1}
    \log\frac{n'(1-q^{1/8})}{q^{1/8}(1-n')}
           =\frac{\ell}{8\lambda}+O_{\mathcal N}(1).
 \]
MP's physical mass budget and phase/response estimates hold
on \([t_s,t_e]\), SRR's weight conversion and late remainder
on the corresponding interval, and RB's core radius and
outside bounds on the entire passage. Hence the other
entry contracts used by CWI have the claimed uniformity.
\end{proof}

\begin{theorem}[Initialized probe error and source rates at entry]
\label{an07psev1:entry}
Define the single center
\[
 \tau_h=t_n+\lambda^{-1}\log\frac{1-n}{n},\qquad
 s_{n'}=\lambda^{-1}\log\frac{n'}{1-n'} .
 \]
Then \(t_{n'}=\tau_h+s_{n'}\) exactly for the auxiliary
logistic family. Uniformly for \(n'\in\mathcal N\),
\(\rho\in[13/20,7/10]\) and \(\zeta\) in WIT's positive
offset interval \([87/50,7/4]\),
\begin{align}
 q^{-2}(E_{\xi_q}-\tfrac12)(t_{n'})
 &=\mathcal E_{\rho,\zeta}(s_{n'})
                             +O_{\mathcal N}(q^{1/20}),\notag\\
 q^{-2}D_p^\tau(t_{n'},\xi_q)
 &=d_{p,\rho,\zeta}(s_{n'})
                             +O_{\mathcal N}(q^{1/20}),\quad p=1,2.
 \label{an07psev1:entrycomparisons}
\end{align}
Rates use the selector convention, or the a.e. interpretation;
no assertion at a gate-hit atom beyond that convention is made.
For common learned witness margins one must still restrict
parameters to WIT's \(\mathcal P_*\) and continue the flow
to its later windows.
\end{theorem}
\begin{proof}
The preceding lemma supplies the hypotheses of CWI
\texttt{corestatic, corerates} uniformly at \(t_{n'}\).
In particular the core mass is \(1+O(q^{11/10})\), the
radius and reference error are \(O_{\mathcal N}(q^{1/20})\),
and current weak weights differ from the fixed incoming law
by \(1+o(1)\). Its selected branch is in the same small
strong box. Since SRR has \(a<28/51\) and \(\delta<a/8\),
the reference gap is larger than
\(87/50-21/34=477/425>1\); the actual core gap is at least
one for small \(q\). SRR's pointwise moment and late
remainder estimates, not an integrated forcing allowance,
are the sourcewise inputs to CWI's core rate calculation.
Thus both core comparisons have normalized error
\(O_{\mathcal N}(q^{1/20})+O(q^2)\).
Apply Theorem \ref{an07psev1:lateprecision}.
Since \(\gamma_*>1/20\),
\[
 q^{-1/20}\epsilon_{\rm bg}(q)
 =q^{8-1/20}+L^2q^{\gamma_*-1/20}\longrightarrow0 .
 \]
This pays the full background without changing the entry rate.
The logistic identity proves the exact common center.
WCS/MP's \(o(1)\) original-clock phase, uniformly on
\(\mathcal N\), is consistent with
\[
 t_{n'}-t_0=Q^{-1}\log\frac{n'}{(1-n')H_0}+o(1).
 \]
Replacing \(Q\) by \(\lambda\) in a fixed-level logarithm
costs \(O_{\mathcal N}(q^2)\); over the long delay the
allowed cost is \(O(q^2\ell)=o(1)\).
No \(O(1)\)-accuracy terminal hitting time is used as a phase.
\end{proof}
Physical time is \((\tau_h+s_{n'})/\mu_1^2\), and the physical
source rate is \(\mu_1^2D_p^\tau\), as in CWI. The result is
uniform only on compact level sets bounded away from zero;
no uniform statement for a new \(q\)-dependent terminal level
is asserted.

\section{Conditional continuation form and S5c handoff}
\label{an07psev1:continuation}
The following statement isolates a reusable inequality. Its
future-time hypotheses are not claimed to be proved here.
\begin{proposition}[Exact conditional energy continuation]
\label{an07psev1:conditional}
On an interval \([a,a+\mathcal T]\) of an exact balanced characteristic
flow, suppose for all relevant receivers
\[
 |R_{12}|\le C_{12}q,\quad |R_{22}^{(1)}|\le C_1q^2,\quad
 |R_{22}^{(2)}|\le C_Ge^{-k_G\sin^2\theta/q^2}
                         \quad\text{when }U_2<0,
 \quad |W_2-U_2|\le d_{\rm al}\sqrt m .
 \]
For every fixed cohort and any \(w>0\),
\begin{align}
 \EE_T(t)\le e^{w(t-a)}\bigg[\EE_T(a)
 +q^2\bigg\{\frac{C_{12}^2}{w}+2C_1
 +\frac{2C_G}{e k_G}
 +\frac{2C_Gd_{\rm al}}{q\sqrt{2e k_G}}\bigg\}
 \int_a^t\mu_T(v)\dd v\bigg].
 \label{an07psev1:conditionalbound}
\end{align}
For fixed \(\mathcal T\), \(d_{\rm al}\le C_{\mathcal T}qL\),
\(w=1\), this is
\[
 \sup_{[a,a+\mathcal T]}\EE_T
 \le e^{\mathcal T}[\EE_T(a)+C_{\mathcal T}q^2L\,\mathcal T
                         \sup_{[a,a+\mathcal T]}\mu_T].
 \]
Thus \(\EE_T(a)=o(q^2)\) and
\(L\sup\mu_T\to0\) suffice for energy continuation;
the CWI output/rate cross terms are paid by
\(L^2\sup\mu_T\to0\) with this alignment scale.
\end{proposition}
\begin{proof}
Repeat the labelwise identity and the Gaussian maxima in
Theorem \ref{an07psev1:energy}, using
\(2C_{12}qh\sqrt m\le wh^2+C_{12}^2q^2m/w\).
Integrate the resulting scalar inequality. This does not
construct a trajectory satisfying the future hypotheses.
\end{proof}

For S5c, the exact sufficient primitive bounds and their
consequences are:
\small
\begin{longtable}{>{\raggedright\arraybackslash}p{.24\textwidth}
 >{\raggedright\arraybackslash}p{.65\textwidth}}
\toprule Quantity and norm & Required scope and consequence\\\midrule\endhead
Cross residuals
& Continue \(\|e_2\|_{L^2(P_1)}+\|e_1\|_{L^2(P_2)}
 \le C_{\mathcal T}q\) on the proposed fixed interval after \(t_b\).
 Loss monotonicity supplies bounded diagonal residuals.
 These imply \(R_{12},R_{21}=O_{\mathcal T}(q)\) for every receiver
 and \(R_{22}^{(1)}=O_{\mathcal T}(q^2)\).\\
Lower-side \(P_2\) moment
& Bounded \(\|e_2\|_{L^2(P_2)}\), the unchanged Gaussian
 model and \(\rho\ge13/20\) imply the exponential matrix
 bound of Lemma \ref{an07psev1:matrix}. It requires no
 output sign, response sign, receiver chart or derivative.\\
Alignment
& Continue \(\sup_b|\theta-\psi|\le C_{\mathcal T}qL\);
 sufficient for the energy proof is
 \(|W_2-U_2|\le C_{\mathcal T}qL\sqrt m\).
 On a fixed interval this follows conditionally from the
 continued cross-residual bound, bounded matrix values,
 and RB's entry alignment: the exact \(\Delta\) equation
 gives only a fixed \(e^{C\mathcal T}\) loss.\\
Fixed-cohort masses
& Require \(L^2\sup\mu_{\Sigma\cup\BB}=o(1)\) for the
 full probe reduction, and the energy seed paid here.
 Under bounded residual matrices, the exact radial equation
 gives \(\mu_T(t)\le e^{C\mathcal T}\mu_T(t_b)\) for any unchanged
 fixed cohort. Thus this mass requirement is a direct
 conditional fixed-time consequence; no new ancestry
 classification or moving-mask flux is hidden in it.\\
Core and primary carriers
& For full-error/background separation, also maintain a
 bounded scaled core chart and
 \(|f_\KK|,|\dot f_{\KK,p}|=O_{\mathcal T}(q)\),
 and primary inactivity, or pay any newly active primary
 output and rates explicitly. These are distinct from the
 lower-half energy inequality.\\
Original-to-leading matching
& Continue masses, reaction/velocity reduction in CWI's
 first-moment norm, selected weak activity and the original
 logit/clock phase on every desired fixed window. This is
 not a consequence of the energy inequality. Source-rate
 transfer is only requested on gate-separated learned
 windows; gate-hit atoms remain excluded.\\\bottomrule
\end{longtable}\normalsize
For the conditional alignment assertion, \(|\operatorname{tr}R|\le C\)
and \(|R_{12}-R_{21}|\le C_{\mathcal T}q\) give
\(D^+|\Delta|\le C|\Delta|+C_{\mathcal T}q\) while \(|\Delta|<1\).
Starting with \(O(qL)\), a fixed interval strictly closes
this comparison for small \(q\). The conditional mass
assertion follows from \(|(\log m)'|\le2\|R\|\).
These are consequences of the listed hypotheses, not a
joint continuation of the original residual guard.

\paragraph{Constants and scope.}
Keep the complete WPT/RB and WCS order unchanged.
The new Gaussian constants \(k_G=169/3200,\ C_G=2\),
layer threshold \(K_{\rm lay}^2=12800/169\), and duration
factor \(31\) are model-only. The pre-entry mass factor
\(e^5\) uses only proved S1 bounds. \(C_\times,C_{\rm a},C_E\)
may depend on the already fixed residual/clock constants;
they multiply \(q^2L^2\) and the proved vanishing masses.
Then fix \(R_\zeta\) and the compact level set \(\mathcal N\),
and choose \(q\) last. Future fixed-window constants may
depend on its fixed duration, not on \(q\).

All requested entry-precision claims close. There is no
unmatched inequality on \([0,t_b]\). The bounds after
\(t_b\) in this section remain conditional S5c obligations.
No later-window original-to-leading theorem, AN8 assembly,
new gate-hit rate transfer or canonical-regime result is asserted.
\end{document}
